\documentclass[a4paper,10pt]{article}
\usepackage[margin=2cm]{geometry}
\usepackage{listings}
\usepackage{xcolor}
\usepackage{fancyhdr}
\usepackage{ctex}
\usepackage{DejaVuSansMono}

% 页眉页脚设置
\pagestyle{fancy}
\fancyhf{}
\fancyhead[C]{\small 审计智鉴流式数据处理系统 V1.0}
\fancyfoot[C]{\thepage/60}
\renewcommand{\headrulewidth}{0.4pt}
\renewcommand{\footrulewidth}{0pt}

% 代码样式设置
\lstset{
    language=Python,
    basicstyle=\ttfamily\footnotesize,
    numbers=left,
    numberstyle=\tiny\color{gray},
    numbersep=8pt,
    frame=none,
    breaklines=true,
    breakatwhitespace=true,
    showstringspaces=false,
    tabsize=4,
    aboveskip=0pt,
    belowskip=0pt,
    lineskip=1pt,
    keywordstyle=\color{blue},
    commentstyle=\color{gray}\itshape,
    stringstyle=\color{orange},
    morekeywords={async,await,True,False,None},
    literate={_}{\_}1
}

\begin{document}
\setcounter{page}{1}

% ==================== backend/main.py ====================
\begin{lstlisting}
"""
FastAPI 主应用入口
财务舞弊识别 SaaS 平台 - Backend API
"""
import sys
import os
# 添加项目根目录到Python路径
sys.path.insert(0, os.path.dirname(
    os.path.dirname(os.path.abspath(__file__))))

from fastapi import FastAPI
from fastapi.middleware.cors import CORSMiddleware
from fastapi.middleware.gzip import GZipMiddleware
from fastapi.staticfiles import StaticFiles
import os

from core.config import settings
from core.database import init_db

# 导入路由
from routers import (user, detection, qa, payment,
    report, user_account, payment_system, membership,
    upload, report_export, financial_statement)

# 创建 FastAPI 应用
app = FastAPI(
    title=settings.APP_NAME,
    version=settings.APP_VERSION,
    description="基于生成式 AI 的上市公司财务舞弊智能识别系统",
    docs_url="/docs",
    redoc_url="/redoc"
)

# 配置 CORS（允许前端访问）
app.add_middleware(
    CORSMiddleware,
    allow_origins=[
        "http://localhost:8501",
        "http://localhost:3000",
        "http://47.76.180.29:8501",
        "http://47.76.180.29:8000",
        "https://*.streamlit.app",
        "https://*.onrender.com"
    ],
    allow_credentials=True,
    allow_methods=["*"],
    allow_headers=["*"],
)

# GZip 压缩中间件
app.add_middleware(GZipMiddleware, minimum_size=1000)

# 挂载静态文件目录
os.makedirs(settings.UPLOAD_FOLDER, exist_ok=True)
os.makedirs("result/reports", exist_ok=True)
os.makedirs("backend/templates/reports", exist_ok=True)
app.mount("/uploads", StaticFiles(
    directory=settings.UPLOAD_FOLDER), name="uploads")
app.mount("/reports", StaticFiles(
    directory="result/reports"), name="reports")

# 包含路由
app.include_router(user.router, prefix=settings.API_PREFIX)
app.include_router(user_account.router, prefix=settings.API_PREFIX)
app.include_router(detection.router, prefix=settings.API_PREFIX)
app.include_router(qa.router, prefix=settings.API_PREFIX)
app.include_router(payment.router, prefix=settings.API_PREFIX)
app.include_router(payment_system.router, prefix=settings.API_PREFIX)
app.include_router(membership.router, prefix=settings.API_PREFIX)
app.include_router(report_export.router, prefix=settings.API_PREFIX)
app.include_router(report.router, prefix=settings.API_PREFIX)
app.include_router(upload.router, prefix=settings.API_PREFIX)
app.include_router(financial_statement.router,
    prefix=settings.API_PREFIX)

# 健康检查
@app.get("/")
def root():
    return {
        "name": settings.APP_NAME,
        "version": settings.APP_VERSION,
        "status": "running"
    }

@app.get("/health")
def health_check():
    return {"status": "healthy"}

# 数据库初始化
@app.on_event("startup")
def startup_event():
    print("=" * 60)
    print(f"{settings.APP_NAME} v{settings.APP_VERSION}")
    print("=" * 60)
    try:
        init_db()
        print("数据库初始化成功")
    except Exception as e:
        print(f"数据库初始化失败：{e}")
    try:
        from services.detection_service import detection_engine
        print("检测引擎初始化成功")
    except Exception as e:
        print(f"检测引擎初始化失败：{e}")

if __name__ == "__main__":
    import uvicorn
    uvicorn.run("main:app", host="0.0.0.0",
        port=8000, reload=settings.DEBUG)
\end{lstlisting}

% ==================== backend/core/security.py ====================
\begin{lstlisting}
"""
用户认证和安全工具
"""
from datetime import datetime, timedelta
from typing import Optional
from jose import JWTError, jwt
from passlib.context import CryptContext
from fastapi import Depends, HTTPException, status
from fastapi.security import HTTPBearer, HTTPAuthorizationCredentials
from sqlalchemy.orm import Session
from backend.core.config import settings
from backend.core.database import get_db
from backend.models.database import User

pwd_context = CryptContext(schemes=["pbkdf2_sha256"],
    deprecated="auto")
security = HTTPBearer()

def verify_password(plain_password: str,
    hashed_password: str) -> bool:
    return pwd_context.verify(plain_password, hashed_password)

def get_password_hash(password: str) -> str:
    return pwd_context.hash(password)

def create_access_token(data: dict,
    expires_delta: Optional[timedelta] = None) -> str:
    to_encode = data.copy()
    if expires_delta:
        expire = datetime.utcnow() + expires_delta
    else:
        expire = datetime.utcnow() + timedelta(
            minutes=settings.ACCESS_TOKEN_EXPIRE_MINUTES)
    to_encode.update({"exp": expire})
    encoded_jwt = jwt.encode(to_encode,
        settings.SECRET_KEY, algorithm=settings.ALGORITHM)
    return encoded_jwt

def decode_access_token(token: str) -> Optional[dict]:
    try:
        payload = jwt.decode(token, settings.SECRET_KEY,
            algorithms=[settings.ALGORITHM])
        return payload
    except JWTError:
        return None

async def get_current_user(
    credentials: HTTPAuthorizationCredentials = Depends(security),
    db: Session = Depends(get_db)
) -> User:
    credentials_exception = HTTPException(
        status_code=status.HTTP_401_UNAUTHORIZED,
        detail="无法验证凭证",
        headers={"WWW-Authenticate": "Bearer"},
    )
    token = credentials.credentials
    payload = decode_access_token(token)
    if payload is None:
        raise credentials_exception
    user_id: str = payload.get("sub")
    if user_id is None:
        raise credentials_exception
    user = db.query(User).filter(User.id == int(user_id)).first()
    if user is None:
        raise credentials_exception
    return user

def get_user_membership_level(user: User) -> str:
    if user.membership_level != "free":
        if user.membership_expire_at and \
                user.membership_expire_at < datetime.utcnow():
            user.membership_level = "free"
            user.membership_expire_at = None
            return "free"
    return user.membership_level

def check_detection_quota(user: User) -> bool:
    membership_level = get_user_membership_level(user)
    if membership_level == "free":
        if user.free_detections_remaining is None or \
                user.free_detections_remaining <= 0:
            return False
        if user.detection_reset_date and \
                user.detection_reset_date < datetime.utcnow().date():
            user.free_detections_remaining = \
                settings.FREE_USER_MONTHLY_DETECTIONS
            user.detection_reset_date = \
                datetime.utcnow().date() + timedelta(days=30)
        return user.free_detections_remaining > 0
    else:
        return True

def consume_detection_quota(user: User) -> int:
    membership_level = get_user_membership_level(user)
    if membership_level == "free":
        if user.free_detections_remaining and \
                user.free_detections_remaining > 0:
            user.free_detections_remaining -= 1
            return user.free_detections_remaining
    return -1

def check_ai_question_quota(user: User) -> bool:
    membership_level = get_user_membership_level(user)
    quotas = {
        "free": settings.FREE_USER_DAILY_AI_QUESTIONS,
        "pro": settings.PRO_USER_DAILY_AI_QUESTIONS,
        "enterprise": settings.ENTERPRISE_USER_DAILY_AI_QUESTIONS,
    }
    quota = quotas.get(membership_level, 0)
    return quota <= 0 or (
        user.free_detections_remaining is None or
        user.free_detections_remaining < quota
    )
\end{lstlisting}

% ==================== backend/models/database.py (部分) ====================
\begin{lstlisting}
"""
数据库模型定义
财务舞弊识别 SaaS 平台 - 核心数据表结构
"""
from sqlalchemy import (Column, Integer, String, Float,
    DateTime, Boolean, Text, ForeignKey, JSON, Date)
from sqlalchemy.dialects.mysql import LONGTEXT
from sqlalchemy.ext.declarative import declarative_base
from sqlalchemy.orm import relationship
from datetime import datetime

Base = declarative_base()

class User(Base):
    __tablename__ = "users"
    id = Column(Integer, primary_key=True, index=True)
    username = Column(String(50), unique=True,
        index=True, nullable=False)
    email = Column(String(100), unique=True, index=True)
    phone = Column(String(20), unique=True, index=True)
    password_hash = Column(String(255), nullable=False)
    user_type = Column(String(20), default="individual")
    membership_level = Column(String(20), default="free")
    membership_expire_at = Column(DateTime)
    balance = Column(Float, default=0.0)
    free_detections_remaining = Column(Integer, default=3)
    detection_reset_date = Column(Date)
    created_at = Column(DateTime, default=datetime.utcnow)
    updated_at = Column(DateTime, default=datetime.utcnow,
        onupdate=datetime.utcnow)
    profile = relationship("UserProfile",
        back_populates="user", uselist=False)
    detections = relationship("DetectionRecord",
        back_populates="user")
    orders = relationship("Order", back_populates="user")
    subscription = relationship("Subscription",
        back_populates="user", uselist=False)
    qa_history = relationship("QAHistory",
        back_populates="user")

class UserProfile(Base):
    __tablename__ = "user_profiles"
    id = Column(Integer, primary_key=True, index=True)
    user_id = Column(Integer, ForeignKey("users.id"),
        unique=True)
    real_name = Column(String(50))
    id_card = Column(String(18))
    id_card_verified = Column(Boolean, default=False)
    company_name = Column(String(200))
    credit_code = Column(String(18))
    company_verified = Column(Boolean, default=False)
    certified = Column(Boolean, default=False)
    certified_at = Column(DateTime)
    created_at = Column(DateTime, default=datetime.utcnow)
    updated_at = Column(DateTime, default=datetime.utcnow,
        onupdate=datetime.utcnow)
    user = relationship("User", back_populates="profile")

class DetectionRecord(Base):
    __tablename__ = "detection_records"
    id = Column(Integer, primary_key=True, index=True)
    user_id = Column(Integer, ForeignKey("users.id"))
    company_name = Column(String(200), nullable=False)
    stock_code = Column(String(10), index=True)
    year = Column(Integer)
    fraud_probability = Column(Float)
    risk_level = Column(String(20))
    risk_score = Column(Float)
    shap_features = Column(JSON)
    ai_feature_scores = Column(JSON)
    risk_labels = Column(JSON)
    risk_evidence_locations = Column(JSON)
    suspicious_segments = Column(JSON)
    ipo_comparison_results = Column(JSON)
    remediation_suggestions = Column(JSON)
    financial_data = Column(JSON)
    mdna_text = Column(LONGTEXT)
    cost = Column(Float, default=0.0)
    detection_type = Column(String(20), default="single")
    status = Column(String(20), default="completed")
    created_at = Column(DateTime, default=datetime.utcnow,
        index=True)
    updated_at = Column(DateTime, default=datetime.utcnow,
        onupdate=datetime.utcnow)
    user = relationship("User", back_populates="detections")
    reports = relationship("Report",
        back_populates="detection")

class Report(Base):
    __tablename__ = "reports"
    id = Column(Integer, primary_key=True, index=True)
    record_id = Column(Integer,
        ForeignKey("detection_records.id"))
    report_type = Column(String(20), default="basic")
    file_path = Column(String(500))
    file_url = Column(String(500))
    is_public = Column(Boolean, default=False)
    share_token = Column(String(64), unique=True)
    share_expire_at = Column(DateTime)
    download_count = Column(Integer, default=0)
    view_count = Column(Integer, default=0)
    created_at = Column(DateTime, default=datetime.utcnow)
    updated_at = Column(DateTime, default=datetime.utcnow,
        onupdate=datetime.utcnow)
    detection = relationship("DetectionRecord",
        back_populates="reports")
\end{lstlisting}

% ==================== backend/routers/user.py ====================
\begin{lstlisting}
"""
用户认证路由
"""
from fastapi import APIRouter, Depends, HTTPException, status
from fastapi.security import HTTPBearer
from sqlalchemy.orm import Session
from datetime import datetime, timedelta

from backend.core.database import get_db
from backend.core.security import (
    verify_password, get_password_hash,
    create_access_token, get_current_user
)
from backend.core.config import settings
from backend.models.database import User, UserProfile
from backend.schemas.schemas import (
    UserCreate, UserLogin, UserResponse, TokenResponse,
    UserUpdate, UserProfileCreate, UserProfileResponse,
    MembershipLevelEnum
)

router = APIRouter(prefix="/user", tags=["用户认证"])
security = HTTPBearer()

@router.post("/register", response_model=UserResponse)
def register(user_data: UserCreate,
    db: Session = Depends(get_db)):
    existing_user = db.query(User).filter(
        User.username == user_data.username).first()
    if existing_user:
        raise HTTPException(
            status_code=status.HTTP_400_BAD_REQUEST,
            detail="用户名已存在"
        )
    if user_data.email:
        existing_email = db.query(User).filter(
            User.email == user_data.email).first()
        if existing_email:
            raise HTTPException(
                status_code=status.HTTP_400_BAD_REQUEST,
                detail="邮箱已被注册"
            )
    if user_data.phone:
        existing_phone = db.query(User).filter(
            User.phone == user_data.phone).first()
        if existing_phone:
            raise HTTPException(
                status_code=status.HTTP_400_BAD_REQUEST,
                detail="手机号已被注册"
            )
    db_user = User(
        username=user_data.username,
        email=user_data.email,
        phone=user_data.phone,
        password_hash=get_password_hash(user_data.password),
        user_type=user_data.user_type.value,
        membership_level="free",
        free_detections_remaining=
            settings.FREE_USER_MONTHLY_DETECTIONS,
        detection_reset_date=datetime.utcnow().date() +
            timedelta(days=30)
    )
    db.add(db_user)
    db.commit()
    db.refresh(db_user)
    db_profile = UserProfile(user_id=db_user.id)
    db.add(db_profile)
    db.commit()
    return db_user

@router.post("/login", response_model=TokenResponse)
def login(login_data: UserLogin,
    db: Session = Depends(get_db)):
    user = db.query(User).filter(
        (User.username == login_data.username) |
        (User.email == login_data.username) |
        (User.phone == login_data.username)
    ).first()
    if not user:
        raise HTTPException(
            status_code=status.HTTP_401_UNAUTHORIZED,
            detail="用户名或密码错误"
        )
    if not verify_password(login_data.password,
            user.password_hash):
        raise HTTPException(
            status_code=status.HTTP_401_UNAUTHORIZED,
            detail="用户名或密码错误"
        )
    access_token = create_access_token(
        data={"sub": str(user.id)}
    )
    return {
        "access_token": access_token,
        "token_type": "bearer",
        "user": user
    }

@router.get("/profile", response_model=UserResponse)
def get_profile(
    current_user: User = Depends(get_current_user),
    db: Session = Depends(get_db)
):
    return current_user

@router.put("/profile", response_model=UserResponse)
def update_profile(
    user_data: UserUpdate,
    current_user: User = Depends(get_current_user),
    db: Session = Depends(get_db)
):
    update_data = user_data.model_dump(exclude_unset=True)
    if user_data.email and user_data.email != \
            current_user.email:
        existing = db.query(User).filter(
            User.email == user_data.email,
            User.id != current_user.id
        ).first()
        if existing:
            raise HTTPException(
                status_code=status.HTTP_400_BAD_REQUEST,
                detail="邮箱已被其他用户使用"
            )
    if user_data.phone and user_data.phone != \
            current_user.phone:
        existing = db.query(User).filter(
            User.phone == user_data.phone,
            User.id != current_user.id
        ).first()
        if existing:
            raise HTTPException(
                status_code=status.HTTP_400_BAD_REQUEST,
                detail="手机号已被其他用户使用"
            )
    for field, value in update_data.items():
        setattr(current_user, field, value)
    db.commit()
    db.refresh(current_user)
    return current_user

@router.post("/logout")
def logout(current_user: User = Depends(get_current_user)):
    return {"message": "登出成功"}
\end{lstlisting}

% ==================== backend/routers/detection.py (核心函数) ====================
\begin{lstlisting}
"""
舞弊检测核心路由
"""
from fastapi import (APIRouter, Depends, HTTPException,
    status, UploadFile, File)
from sqlalchemy.orm import Session
from typing import Optional, List, Dict, Any, Tuple
from functools import lru_cache
import json
import re

from backend.core.database import get_db
from backend.core.security import get_current_user
from backend.core.config import settings
from backend.models.database import (User, DetectionRecord,
    DemoCase, Report)
from backend.schemas.schemas import (
    DetectionCreate, DetectionResponse,
    DetectionDetailResponse, RiskLevelEnum
)

_cases_cache = {}
_cases_cache_time = 0
import time

def get_cached_cases(db: Session, featured_only: bool = False):
    global _cases_cache, _cases_cache_time
    cache_key = f"featured_{featured_only}"
    current_time = time.time()
    if cache_key in _cases_cache and \
            (current_time - _cases_cache_time) < 300:
        return _cases_cache[cache_key]
    query = db.query(DemoCase)
    if featured_only:
        query = query.filter(DemoCase.is_featured == True)
    cases = query.order_by(DemoCase.sort_order,
        DemoCase.id).all()
    _cases_cache[cache_key] = cases
    _cases_cache_time = current_time
    return cases

router = APIRouter(prefix="/detection", tags=["舞弊检测"])

def compute_shap_features(ai_scores: dict,
    financial_data: dict) -> dict:
    shap_values = {}
    for feature, score in ai_scores.items():
        if isinstance(score, (list, tuple)):
            score = score[0]
        if isinstance(score, str):
            safe_score = 0.0
        else:
            try:
                safe_score = float(score)
            except:
                safe_score = 0.0
        shap_values[feature] = round(safe_score * 0.3, 4)
    if financial_data:
        cash = float(financial_data.get("货币资金", 0) or 0)
        short_loan = float(
            financial_data.get("短期借款", 0) or 0)
        if cash > 50 and short_loan > 20:
            shap_values["存贷双高指标"] = 0.25
        net_cash = float(financial_data.get(
            "经营活动现金流净额", 0) or 0)
        net_profit = float(financial_data.get(
            "净利润", 0) or 0)
        if net_profit > 0 and net_cash < 0:
            shap_values["现金流背离指标"] = 0.22
        inventory = float(financial_data.get("存货", 0) or 0)
        total_assets = float(financial_data.get(
            "总资产", 1) or 1)
        if total_assets > 0 and inventory / total_assets > 0.15:
            shap_values["存货占比指标"] = 0.18
    total = sum(shap_values.values())
    if total > 0:
        shap_values = {k: round(v / total, 4)
            for k, v in shap_values.items()}
    return shap_values

async def extract_ai_features_from_text(
    mdna_text: str, financial_data: dict) -> dict:
    from backend.services.detection_service import \
        detection_engine
    if not mdna_text:
        return detection_engine._get_default_ai_features()
    try:
        ai_features = await detection_engine.extract_ai_features(
            mdna_text, financial_data)
        return ai_features
    except Exception as e:
        print(f"LLM特征提取失败: {e}")
        return _fallback_ai_feature_extraction(
            mdna_text, financial_data)

def _fallback_ai_feature_extraction(
    mdna_text: str, financial_data: dict) -> dict:
    ai_scores = {
        "CON_SEM_AI": 0.35,
        "COV_RISK_AI": 0.35,
        "TONE_ABN_AI": 0.35,
        "FIT_TD_AI": 0.35,
        "HIDE_REL_AI": 0.35,
        "DEN_ABN_AI": 0.35,
        "STR_EVA_AI": 0.35
    }
    text_lower = mdna_text.lower()
    contradiction_keywords = ["但是", "然而", "尽管",
        "虽然", "不过", "却", "反之"]
    contra_count = sum(1 for kw in contradiction_keywords
        if kw in text_lower)
    if contra_count >= 2:
        ai_scores["CON_SEM_AI"] = min(
            0.4 + contra_count * 0.08, 0.9)
    risk_keywords = ["风险", "不确定性", "挑战",
        "困难", "压力", "波动", "下滑"]
    risk_count = sum(1 for kw in risk_keywords
        if kw in text_lower)
    if risk_count < 3:
        ai_scores["COV_RISK_AI"] = 0.65
    elif risk_count > 8:
        ai_scores["COV_RISK_AI"] = 0.55
    positive_keywords = ["大幅增长", "突破", "领先",
        "优异", "显著", "持续向好", "创新高"]
    positive_count = sum(1 for kw in positive_keywords
        if kw in text_lower)
    if positive_count > 3:
        ai_scores["TONE_ABN_AI"] = min(
            0.45 + positive_count * 0.1, 0.95)
    if financial_data:
        revenue_growth = float(financial_data.get(
            "营业收入增长率", 0) or 0)
        net_profit = float(financial_data.get(
            "净利润", 0) or 0)
        cash_flow = float(financial_data.get(
            "经营活动现金流净额", 0) or 0)
        if ("增长" in text_lower or "提升" in text_lower) \
                and revenue_growth < -0.1:
            ai_scores["FIT_TD_AI"] = 0.85
        if net_profit > 0 and cash_flow < 0:
            ai_scores["FIT_TD_AI"] = min(
                ai_scores["FIT_TD_AI"] + 0.25, 0.9)
    related_keywords = ["关联方", "关联交易", "关联关系",
        "少数股东", "实际控制人", "控股股东"]
    related_count = sum(1 for kw in related_keywords
        if kw in text_lower)
    if related_count == 0:
        ai_scores["HIDE_REL_AI"] = 0.55
    elif related_count > 5:
        ai_scores["HIDE_REL_AI"] = min(
            0.45 + related_count * 0.06, 0.8)
    text_length = len(mdna_text)
    if text_length < 300:
        ai_scores["DEN_ABN_AI"] = 0.6
    elif text_length > 3000:
        ai_scores["DEN_ABN_AI"] = 0.5
    evade_keywords = ["可能", "或许", "大概", "预计",
        "计划", "拟", "将", "有望"]
    evade_count = sum(1 for kw in evade_keywords
        if kw in text_lower)
    if evade_count > 8:
        ai_scores["STR_EVA_AI"] = min(
            0.4 + evade_count * 0.05, 0.85)
    return ai_scores

@router.get("/cases", response_model=List[DemoCaseResponse])
def get_demo_cases(
    featured_only: bool = False,
    db: Session = Depends(get_db)
):
    return get_cached_cases(db, featured_only)

@router.get("/cases/{case_id}",
    response_model=DemoCaseDetailResponse)
def get_demo_case(case_id: int,
    db: Session = Depends(get_db)):
    case = db.query(DemoCase).filter(
        DemoCase.id == case_id).first()
    if not case:
        raise HTTPException(
            status_code=status.HTTP_404_NOT_FOUND,
            detail="案例不存在"
        )
    return case
\end{lstlisting}

% ==================== backend/services/detection_service.py (核心类) ====================
\begin{lstlisting}
"""
舞弊检测核心功能实现
包含 AI 特征提取、风险评分计算、SHAP 分析等
"""
import asyncio
import json
import re
import os
import hashlib
from typing import Dict, List, Optional, Tuple, Any
from functools import lru_cache
import numpy as np
import shap
import xgboost as xgb
import joblib
from sqlalchemy.orm import Session
from cachetools import TTLCache

from backend.core.config import settings
from backend.models.database import DetectionRecord, DemoCase
from backend.core.cache_manager import (get_llm_cache,
    get_shap_cache, cached)

ai_result_cache = get_llm_cache()
shap_cache = get_shap_cache()
from backend.schemas.schemas import DetectionCreate, RiskLevelEnum
from backend.services.analysis_service import analysis_service
from backend.services.detailed_shap_analysis import \
    get_detailed_shap_analysis

PROJECT_ROOT = os.path.dirname(os.path.dirname(
    os.path.dirname(os.path.abspath(__file__))))
MODEL_DIR = os.path.join(PROJECT_ROOT, "result", "models")

if not os.path.exists(MODEL_DIR):
    alt_model_dir = os.path.join(PROJECT_ROOT, "..", "models")
    if os.path.exists(alt_model_dir):
        MODEL_DIR = os.path.abspath(alt_model_dir)

class FraudDetectionEngine:
    def __init__(self):
        self.ai_model = None
        self.traditional_model = None
        self.feature_names = []
        self.scaler = None
        self.selected_features = None
        self.numeric_columns = None
        self._load_models()

    def _load_models(self):
        try:
            ai_model_path = os.path.join(MODEL_DIR,
                "model_ai_XGBoost.pkl")
            trad_model_path = os.path.join(MODEL_DIR,
                "model_trad_XGBoost.pkl")
            scaler_path = os.path.join(MODEL_DIR, "scaler.pkl")
            features_path = os.path.join(MODEL_DIR,
                "selected_features.pkl")
            numeric_path = os.path.join(MODEL_DIR,
                "numeric_columns.pkl")
            if os.path.exists(ai_model_path):
                self.ai_model = joblib.load(ai_model_path)
            if os.path.exists(trad_model_path):
                self.traditional_model = joblib.load(
                    trad_model_path)
            if os.path.exists(scaler_path):
                self.scaler = joblib.load(scaler_path)
            if os.path.exists(features_path):
                self.selected_features = joblib.load(
                    features_path)
            if os.path.exists(numeric_path):
                self.numeric_columns = joblib.load(
                    numeric_path)
            if self.ai_model:
                try:
                    self.feature_names = \
                        self.ai_model.get_booster().feature_names
                except Exception:
                    try:
                        self.feature_names = \
                            self.ai_model.feature_names_in_
                    except:
                        self.feature_names = []
        except Exception as e:
            self.ai_model = xgb.XGBClassifier()
            self.traditional_model = xgb.XGBClassifier()

    @cached(cache_name="llm_results", maxsize=200,
            ttl=86400)
    async def extract_ai_features(self, mdna_text: str,
        financial_data: Dict[str, Any]) -> Dict[str, float]:
        if not mdna_text:
            return self._get_default_ai_features()
        cache_key = self._generate_cache_key(
            mdna_text, financial_data)
        if cache_key in ai_result_cache:
            return ai_result_cache[cache_key]
        prompt = settings.OPTIMIZED_PROMPT_TEMPLATE.format(
            mdna_text=mdna_text[:1000],
            financial_data=json.dumps(financial_data,
                ensure_ascii=False)
        )
        try:
            response = await self._call_llm_api(prompt)
            features = self._parse_llm_response(response)
            ai_result_cache[cache_key] = features
            return features
        except Exception as e:
            return self._get_default_ai_features()

    def _generate_cache_key(self, mdna_text: str,
        financial_data: Dict[str, Any]) -> str:
        text_prefix = mdna_text[:200].strip()
        key_data = f"{text_prefix}_{json.dumps(
            financial_data, sort_keys=True,
            ensure_ascii=False)}"
        return hashlib.md5(key_data.encode()).hexdigest()

    async def _call_llm_api(self, prompt: str) -> str:
        import httpx
        try:
            async with httpx.AsyncClient(
                    timeout=60.0) as client:
                response = await client.post(
                    f"{settings.DASHSCOPE_BASE_URL}/chat/completions",
                    headers={
                        "Authorization": f"Bearer {settings.DASHSCOPE_API_KEY}",
                        "Content-Type": "application/json"
                    },
                    json={
                        "model": settings.MODEL_QWEN_DETECTION,
                        "messages": [
                            {"role": "system", "content":
                                "你是财务舞弊识别领域的专家，擅长通过文本分析识别财务风险信号。"},
                            {"role": "user", "content": prompt}
                        ],
                        "temperature": 0.3,
                        "max_tokens": 800
                    }
                )
                if response.status_code == 200:
                    result = response.json()
                    return result["choices"][0]["message"]["content"]
                else:
                    return self._get_fallback_ai_response()
        except Exception as e:
            return self._get_fallback_ai_response()

    def _get_fallback_ai_response(self) -> str:
        return json.dumps({
            "CON_SEM_AI": 0.45,
            "COV_RISK_AI": 0.42,
            "TONE_ABN_AI": 0.40,
            "FIT_TD_AI": 0.44,
            "HIDE_REL_AI": 0.41,
            "DEN_ABN_AI": 0.43,
            "STR_EVA_AI": 0.42,
            "analysis_notes": "API调用失败，使用默认评估值"
        })

    def _parse_llm_response(self, response: str) -> Dict[str, float]:
        cleaned_response = re.sub(r'^```json\n|\n```$',
            '', response.strip())
        features = json.loads(cleaned_response)
        standardized = {}
        for feature in settings.WEIGHTED_FEATURES.keys():
            score = features.get(feature, 0.3)
            standardized[feature] = min(max(
                float(score), 0.0), 1.0)
        if 'key_risks' in features:
            standardized['_key_risks'] = features['key_risks']
        if 'text_evidence' in features:
            standardized['_text_evidence'] = \
                features['text_evidence']
        return standardized

    def _get_default_ai_features(self) -> Dict[str, float]:
        return {feature: 0.3
            for feature in settings.WEIGHTED_FEATURES.keys()}

    def calculate_fraud_probability(self,
        financial_data: Dict[str, Any],
        ai_features: Dict[str, float]
    ) -> Tuple[float, RiskLevelEnum, List[Dict[str, Any]], float]:
        trad_score = self._calculate_traditional_risk(
            financial_data)
        ai_score = 0.0
        for feature, score in ai_features.items():
            weight = settings.WEIGHTED_FEATURES.get(feature, 1.0)
            ai_score += score * weight
        ai_score = ai_score / sum(
            settings.WEIGHTED_FEATURES.values()) * 65
        base_risk = 35
        total_score = base_risk + trad_score + ai_score
        total_score = min(total_score, 100)
        normalized_score = total_score / 100
        fraud_probability = min(max(
            0.50 + normalized_score * 0.5, 0), 0.95)
        if fraud_probability >= 0.55:
            risk_level = RiskLevelEnum.HIGH
        elif fraud_probability >= 0.30:
            risk_level = RiskLevelEnum.MEDIUM
        else:
            risk_level = RiskLevelEnum.LOW
        risk_labels = self._generate_risk_labels(
            financial_data, ai_features, fraud_probability)
        return (fraud_probability, risk_level,
            risk_labels, total_score)

    def _calculate_traditional_risk(self,
        financial_data: Dict[str, Any]) -> float:
        score = 0.0
        cash = float(financial_data.get("货币资金", 0) or 0)
        short_loan = float(financial_data.get(
            "短期借款", 0) or 0)
        if cash > 50 and short_loan > 20:
            score += 20
        net_cash = float(financial_data.get(
            "经营活动现金流净额", 0) or 0)
        net_profit = float(financial_data.get(
            "净利润", 0) or 0)
        if net_profit > 0 and net_cash < 0:
            score += 15
        inventory = float(financial_data.get("存货", 0) or 0)
        total_assets = float(financial_data.get(
            "总资产", 1) or 1)
        if total_assets > 0 and inventory / total_assets > 0.15:
            score += 10
        return min(score, 40)

    def _generate_risk_labels(self, financial_data, ai_features,
        fraud_probability):
        risk_labels = []
        cash = float(financial_data.get("货币资金", 0) or 0)
        short_loan = float(financial_data.get(
            "短期借款", 0) or 0)
        net_cash = float(financial_data.get(
            "经营活动现金流净额", 0) or 0)
        net_profit = float(financial_data.get(
            "净利润", 0) or 0)
        if cash > 50 and short_loan > 20:
            risk_labels.append({
                "label": "存贷双高",
                "score": 0.85,
                "description": "货币资金和短期借款同时处于高位"
            })
        if net_profit > 0 and net_cash < 0:
            risk_labels.append({
                "label": "现金流背离",
                "score": 0.90,
                "description": "净利润为正但经营现金流为负"
            })
        label_map = {
            "CON_SEM_AI": ("文本语义矛盾",
                "MD&A 文本中存在前后矛盾的表述"),
            "FIT_TD_AI": ("文本-数据不一致",
                "文本描述与财务数据不匹配"),
            "COV_RISK_AI": ("风险披露不足",
                "未充分披露重大风险因素"),
            "HIDE_REL_AI": ("关联交易隐藏",
                "关联交易披露不充分或存在隐藏"),
            "TONE_ABN_AI": ("语调异常乐观",
                "管理层语调过于乐观"),
            "DEN_ABN_AI": ("信息密度异常",
                "信息披露过于简略或冗长"),
            "STR_EVA_AI": ("回避表述",
                "对关键问题使用模糊、回避性表述")
        }
        thresholds = {
            "CON_SEM_AI": 0.6, "COV_RISK_AI": 0.6,
            "TONE_ABN_AI": 0.55, "FIT_TD_AI": 0.6,
            "HIDE_REL_AI": 0.55, "DEN_ABN_AI": 0.65,
            "STR_EVA_AI": 0.6
        }
        for feature, score in ai_features.items():
            if feature.startswith('_') or \
                    feature not in label_map:
                continue
            try:
                score_val = float(score) if score is not None \
                    else 0.0
            except:
                score_val = 0.0
            if score_val >= thresholds.get(feature, 0.6):
                label, desc = label_map[feature]
                risk_labels.append({
                    "label": label, "score": round(score_val, 2),
                    "description": desc
                })
        return risk_labels[:10]
\end{lstlisting}

% ==================== backend/routers/qa.py ====================
\begin{lstlisting}
"""
AI 智能问答路由 - 支持流式输出
"""
from fastapi import APIRouter, Depends, HTTPException, status
from fastapi.responses import StreamingResponse
from sqlalchemy.orm import Session
from typing import List, Optional, AsyncGenerator
import httpx
import json

from backend.core.database import get_db
from backend.core.security import get_current_user
from backend.core.config import settings
from backend.models.database import User, QAHistory
from backend.schemas.schemas import (QAAskRequest, QAResponse,
    QASuggestionResponse)

router = APIRouter(prefix="/qa", tags=["AI 问答"])

PRESET_QUESTIONS = {
    "theory": [
        "什么是财务舞弊的信号传递理论？",
        "信息不对称如何导致财务舞弊？",
        "舞弊三角理论的核心要素是什么？",
        "GONE 理论如何解释财务舞弊？"
    ],
    "practice": [
        "什么是存贷双高？如何识别？",
        "如何识别 MD&A 中的语义矛盾？",
        "审计中如何使用 AI 工具辅助？",
        "如何通过现金流分析识别舞弊？"
    ],
    "policy": [
        "A 股信息披露规则有哪些核心要求？",
        "证监会财务舞弊处罚标准是什么？",
        "新证券法对舞弊的处罚力度如何？",
        "会计师事务所的审计责任有哪些？"
    ],
    "case": [
        "康美药业舞弊案的关键识别点是什么？",
        "瑞幸咖啡虚增收入的手法是什么？",
        "獐子岛扇贝跑路事件如何识别？",
        "贵州茅台为什么是健康企业典型？"
    ],
    "platform": [
        "如何解读舞弊概率？",
        "SHAP 分析的原理是什么？",
        "如何导出检测报告？",
        "会员版和免费版有什么区别？"
    ]
}

SYSTEM_PROMPT = """你是一位财务舞弊识别领域的专家助手。
请用专业、准确但通俗易懂的语言回答用户问题。"""

async def call_llm_api(question: str, context: str = "",
    image_base64: str = None, image_mime: str = None) -> str:
    user_prompt = f"{context}\n\n问题：{question}" \
        if context else question
    is_multimodal = image_base64 is not None and \
        image_mime is not None
    model = settings.MODEL_QWEN_VL if is_multimodal \
        else settings.MODEL_QWEN_FS
    if is_multimodal:
        messages = [
            {"role": "system", "content": SYSTEM_PROMPT},
            {"role": "user", "content": [
                {"type": "text", "text": user_prompt},
                {"type": "image_url", "image_url": {
                    "url": f"data:{image_mime};base64,{image_base64}"
                }}
            ]}
        ]
    else:
        messages = [
            {"role": "system", "content": SYSTEM_PROMPT},
            {"role": "user", "content": user_prompt}
        ]
    try:
        async with httpx.AsyncClient(timeout=60.0) as client:
            response = await client.post(
                f"{settings.DASHSCOPE_BASE_URL}/chat/completions",
                headers={
                    "Authorization": f"Bearer {settings.DASHSCOPE_API_KEY}",
                    "Content-Type": "application/json"
                },
                json={
                    "model": model,
                    "messages": messages,
                    "temperature": 0.7,
                    "max_tokens": 1500
                }
            )
            if response.status_code == 200:
                result = response.json()
                return result["choices"][0]["message"]["content"]
            else:
                return "抱歉，暂时无法回答该问题。"
    except Exception as e:
        return f"抱歉，回答失败：{str(e)[:100]}"

async def call_llm_api_streaming(question: str,
    context: str = "") -> AsyncGenerator[str, None]:
    user_prompt = f"{context}\n\n问题：{question}" \
        if context else question
    try:
        async with httpx.AsyncClient(timeout=120.0) as client:
            async with client.stream("POST",
                f"{settings.DASHSCOPE_BASE_URL}/chat/completions",
                headers={
                    "Authorization": f"Bearer {settings.DASHSCOPE_API_KEY}",
                    "Content-Type": "application/json"
                },
                json={
                    "model": settings.MODEL_QWEN_FS,
                    "messages": [
                        {"role": "system",
                            "content": SYSTEM_PROMPT},
                        {"role": "user", "content": user_prompt}
                    ],
                    "temperature": 0.7,
                    "max_tokens": 1500,
                    "stream": True
                }
            ) as response:
                if response.status_code != 200:
                    yield f"data: {json.dumps({'error':
                        'API调用失败'})}\n\n"
                    return
                async for line in response.aiter_lines():
                    if not line or line.strip() == "":
                        continue
                    if line.startswith("data: "):
                        data_str = line[6:]
                        if data_str == "[DONE]":
                            yield f"data: {json.dumps(
                                {'done': True})}\n\n"
                            break
                        try:
                            data = json.loads(data_str)
                            if "choices" in data and \
                                    len(data["choices"]) > 0:
                                delta = data["choices"][0].get(
                                    "delta", {})
                                content = delta.get("content")
                                if content and \
                                        isinstance(content, str):
                                    yield f"data: {json.dumps(
                                        {'content': content})}\n\n"
                        except json.JSONDecodeError:
                            continue
    except Exception as e:
        yield f"data: {json.dumps({'error': str(e)})}\n\n"

@router.post("/ask", response_model=QAResponse)
async def ask_question(
    question_data: QAAskRequest,
    current_user: User = Depends(get_current_user),
    db: Session = Depends(get_db)
):
    if question_data.category and question_data.category in \
            ["theory", "practice", "policy", "case", "platform"]:
        context = f"问题类别：{question_data.category}"
    else:
        context = ""
    answer = await call_llm_api(
        question_data.question, context,
        image_base64=question_data.image_base64,
        image_mime=question_data.image_mime
    )
    db_qa = QAHistory(
        user_id=current_user.id,
        question=question_data.question,
        answer=answer,
        category=question_data.category
    )
    db.add(db_qa)
    db.commit()
    db.refresh(db_qa)
    return db_qa

@router.post("/ask-stream")
async def ask_question_stream(
    question_data: QAAskRequest,
    current_user: User = Depends(get_current_user),
    db: Session = Depends(get_db)
):
    if question_data.category and question_data.category in \
            ["theory", "practice", "policy", "case", "platform"]:
        context = f"问题类别：{question_data.category}"
    else:
        context = ""
    async def event_generator() -> AsyncGenerator[str, None]:
        full_answer = ""
        yield f"data: {json.dumps({'start': True})}\n\n"
        async for chunk in call_llm_api_streaming(
                question_data.question, context):
            yield chunk
            try:
                if chunk.startswith("data: "):
                    data = json.loads(chunk[6:])
                    content = data.get("content")
                    if content and isinstance(content, str):
                        full_answer += content
            except:
                pass
        yield f"data: {json.dumps({'done': True})}\n\n"
        try:
            db_qa = QAHistory(
                user_id=current_user.id,
                question=question_data.question,
                answer=full_answer or "（流式输出内容）",
                category=question_data.category
            )
            db.add(db_qa)
            db.commit()
        except Exception as e:
            print(f"保存问答历史失败: {e}")
    return StreamingResponse(
        event_generator(),
        media_type="text/event-stream",
        headers={
            "Cache-Control": "no-cache",
            "Connection": "keep-alive",
            "X-Accel-Buffering": "no"
        }
    )

@router.get("/history", response_model=List[QAResponse])
def get_qa_history(
    page: int = 1, page_size: int = 20,
    current_user: User = Depends(get_current_user),
    db: Session = Depends(get_db)
):
    query = db.query(QAHistory).filter(
        QAHistory.user_id == current_user.id
    ).order_by(QAHistory.created_at.desc())
    offset = (page - 1) * page_size
    return query.offset(offset).limit(page_size).all()

@router.get("/suggestions",
    response_model=List[QASuggestionResponse])
def get_suggestions(
    category: Optional[str] = None,
    current_user: User = Depends(get_current_user)
):
    if category and category in PRESET_QUESTIONS:
        return [QASuggestionResponse(
            category=category,
            questions=PRESET_QUESTIONS[category]
        )]
    return [QASuggestionResponse(category=cat,
        questions=questions)
        for cat, questions in PRESET_QUESTIONS.items()]
\end{lstlisting}

\end{document}
